% ============================================================
% Chapter 8: Sustainability and Ethical Implications
%   Section: Sustainability Report
%     Subsection: Social Dimension
% Included from chapters/08-sustainability-ethics/01-sustainability-report/sustainability-report.tex
% ============================================================

\subsection{Social Dimension}

\paragraph{Who benefits.} Readers who want news that is constructive
and verified get articles whose claims carry a verdict and the evidence
behind it, instead of a promise. Journalists and small newsrooms get a
tool that reads the sources, screens what does not fit, proposes the
claims worth checking and does a first search for evidence --- the
steps that take most of the time before any judgement is made.

\paragraph{Who could be harmed.} Three groups. Readers, if a wrong
verdict is presented with more authority than it deserves; the system
answers by preferring ``not verified'', by setting failures apart and
by showing every verdict's evidence, but the X-Fact pilot shows that
the current verifier is often wrong, and a reader shown a confident
\texttt{TRUE} cannot tell. Journalists, if a tool like this is used to
reduce newsrooms rather than to support them; RQ4 is asked as a
question of extent precisely because the results show where a person is
still needed. And the sources, whose articles are read automatically;
how that is limited is in Section~\ref{sec:ethics}.

\paragraph{Trust.} A verdict that can be inspected is the point of the
system: a reader can follow each claim to its sources and see why each
was trusted, including when a source's reliability is unknown. Whether
that transparency increases readers' trust, or only informs those who
already look, has not been studied here.
